%%=============================================================================
%% Conclusie
%%=============================================================================

\chapter{Conclusie}\label{ch:conclusie}

In deze bachelorproef werd de ontwikkeling van een \acrfull{gis}-gebaseerde tool ter ondersteuning van de energieanalyse voor de \acrfull{et} kandidaat locatie in de \acrfull{emr} beschreven. Het analyseren van grensoverschrijdende datasets uit verschillende bronnen, met uiteenlopende formaten en kwaliteitsniveaus, vormde hierbij een uitdaging. Deze uitdaging werd vergroot door de noodzaak om snel meerdere opties voor duurzame energiewinning te analyseren.\bigskip

Uit de literatuurstudie en vereistenanalyse bleek dat een duurzame, trillingsvrije en beschikbare energievoorziening essentieel is voor de werking van de \acrshort{et}. Hierdoor werden bepaalde energiebronnen, zoals nucleaire- en windenergie, minder geschikt bevonden binnen de context van dit onderzoek. Zonne-energie kwam naar voren als een veelbelovende optie vanwege de flexibiliteit in de plaatsing en configuratie van installaties, de schaalbaarheid en beperkte impact op de omgeving. Daarnaast maakt deze technologie dubbel landgebruik mogelijk, waarbij energieproductie wordt gecombineerd met andere ruimtelijke functies op dezelfde locatie.\bigskip

Op basis van een vergelijking van bestaande \acrshort{gis}-tools afgetoetst aan de noden van de \acrshort{et}-onderzoeksgroep werd QGIS gekozen als fundament, aangevuld met de ingebouwde PyQGIS Python bibliotheek, voor de ontwikkeling van aangepaste functionaliteiten. Deze keuze maakte het mogelijk om een tool te ontwikkelen die zowel krachtig als toegankelijk is. De implementatie resulteerde in een \acrfull{poc} tool waarmee gebruikers geografische datasets kunnen analyseren binnen een gekozen of gegenereerd zoekgebied. De resultaten kunnen worden gevisualiseerd en geëxporteerd naar gangbare en open formaten zoals CSV en PDF.\bigskip

Een belangrijke meerwaarde van de ontwikkelde tool is de gebruiksvriendelijke interface, die ook gebruikers zonder specifieke kennis van \acrshort{gis} of programmeren in staat stelt om analyses uit te voeren. Daarnaast werd veel aandacht besteed aan modulariteit en uitbreidbaarheid, waardoor de tool eenvoudig aangepast of uitgebreid kan worden aan toekomstige noden en inzichten.\bigskip

Desondanks kent de bekomen tool ook enkele beperkingen. De tool richt zich momenteel voornamelijk op oppervlakte analyse van geografische data en bevat nog geen uitgewerkte functionaliteiten voor het berekenen van energiepotentieel. Bovendien blijft de kwaliteit van de resultaten afhankelijk van de beschikbaarheid en nauwkeurigheid van de gebruikte datasets. Ook moeten de onderzoekers nog manueel waken over mogelijke overlap van potentiële locaties.\bigskip

Samenvattend toont deze thesis aan dat een \acrshort{gis}-gebaseerde analysetool een waardevolle bijdrage kan leveren aan het energieplanningsproces voor de \acrshort{et}. Toekomstig onderzoek kan zich richten op het integreren van geavanceerde analysemethoden, zoals \acrfull{mcda} en energieproductiemodellen, alsook op het harmoniseren en uniformeren van gelijkaardige datasets uit verschillende bronnen, om de tool verder te versterken en te verfijnen.